\usepackage{amsmath,array,xparse}

\ExplSyntaxOn
\tl_new:N \l__linsys_body_tl
\cs_new:Npn \sysequals { = }

% Write ordinary = signs and separate equations with \\.
% The environment adds a dotted row at every line break.
% Do not add a trailing \\ after the last equation.
\NewDocumentEnvironment{sys}{b}
  {
    \tl_set:Nn \l__linsys_body_tl {#1}
    \tl_replace_all:Nnn \l__linsys_body_tl { = } { & {}={} & }
    \tl_replace_all:Nnn \l__linsys_body_tl { \\ }
      { \\ & \vdots & \\ }
    % Let \lineq also align correctly inside this environment.
    \cs_set:Npn \sysequals { & {}={} & }
    \left\{
    \begin{array}{@{}>{\displaystyle}r@{}c@{}>{\displaystyle}l@{}}
      \tl_use:N \l__linsys_body_tl
    \end{array}
    \right.
  }
  {}
\ExplSyntaxOff

% Generic m-by-n system: coefficients, unknowns, right-hand side.
\NewDocumentCommand{\linsys}{G{a} G{x} G{b} G{m} G{n}}{%
  \begingroup
  \IfBlankTF{#1}{\def\linsysA{a}}{\def\linsysA{#1}}%
  \IfBlankTF{#2}{\def\linsysX{x}}{\def\linsysX{#2}}%
  \IfBlankTF{#3}{\def\linsysB{b}}{\def\linsysB{#3}}%
  \IfBlankTF{#4}{\def\linsysM{m}}{\def\linsysM{#4}}%
  \IfBlankTF{#5}{\def\linsysN{n}}{\def\linsysN{#5}}%

  \def\linsysZero{0}%
  \ifx\linsysB\linsysZero
    \def\linsysFirstRHS{0}%
    \def\linsysLastRHS{0}%
  \else
    \def\linsysFirstRHS{{\linsysB}_1}%
    \def\linsysLastRHS{{\linsysB}_{\linsysM}}%
  \fi

  \begin{sys}
    {\linsysA}_{11}{\linsysX}_1+\cdots
      +{\linsysA}_{1\linsysN}{\linsysX}_{\linsysN}
      =\linsysFirstRHS\\
    {\linsysA}_{\linsysM 1}{\linsysX}_1+\cdots
      +{\linsysA}_{\linsysM\linsysN}{\linsysX}_{\linsysN}
      =\linsysLastRHS
  \end{sys}%
  \endgroup
}

% Generic vertical vector
\NewDocumentCommand{\hvc}{m G{n}}{%
  \left({#1}_1,\ldots,{#1}_{\IfBlankTF{#2}{n}{#2}}\right)%
}

\NewDocumentCommand{\vvc}{m G{n}}{%
  \begin{pmatrix}
    {#1}_1\\
    \vdots\\
    {#1}_{\IfBlankTF{#2}{n}{#2}}
  \end{pmatrix}%
}

% Put your own mathematical commands here.
\newcommand{\R}{\mathbb{R}}
\newcommand{\C}{\mathbb{C}}
